\chapter{Error-Correcting Codes — Core Definitions}
%%%%%%%%%%%%%%%%%%%%%%%%%%%%%%%%%%%%%%%%%%%%%%%%%%%%%%%%%%%%%%%%%%%%%%%%%%%%%%%
\section{Overview}

This file documents the foundational definitions for $q$-ary block codes over a finite field
alphabet $\alpha$.  Throughout, $\alpha$ is a type carrying instances
\lean{[Fintype α]}, \lean{[Field α]}, \lean{[DecidableEq α]}, and \lean{[Nonempty α]}.

%%%%%%%%%%%%%%%%%%%%%%%%%%%%%%%%%%%%%%%%%%%%%%%%%%%%%%%%%%%%%%%%%%%%%%%%%%%%%%%
\section{Basic objects: codewords and codes}

\subsection{Codewords}

\begin{definition}[Codeword]
\lean{ErrorCorrectingCodes.Codeword}
\leanok
A codeword of length $n$ over alphabet $\alpha$ is a function $c : \mathrm{Fin}\,n \to \alpha$.
\end{definition}

\begin{definition}[Pointwise operations]
\lean{ErrorCorrectingCodes.Codeword.add}
\lean{ErrorCorrectingCodes.Codeword.sub}
\lean{ErrorCorrectingCodes.Codeword.zero}
\leanok
Pointwise addition, subtraction, and the all-zero codeword.
\end{definition}

\subsection{Codes and linear codes}

\begin{definition}[Code]
\lean{ErrorCorrectingCodes.Codeword.Code}
\leanok
A (block) code of length $n$ over $\alpha$ is a finite set (\lean{Finset}) of codewords.
\end{definition}

\begin{definition}[Linear code via generator matrix]
\lean{ErrorCorrectingCodes.Codeword.Linear_Code}
\leanok
A code $C$ is linear with generator matrix $G$ (shape $n\times m$) if every message
$c':\mathrm{Fin}\,m\to\alpha$ maps to a codeword $Gc'\in C$, and every codeword in $C$
arises this way.
\end{definition}

\begin{definition}[Linear code, existential form]
\lean{ErrorCorrectingCodes.Codeword.Linear_Code'}
\leanok
Existential version: there exists some $m$ and generator matrix $G$ witnessing linearity.
\end{definition}

%%%%%%%%%%%%%%%%%%%%%%%%%%%%%%%%%%%%%%%%%%%%%%%%%%%%%%%%%%%%%%%%%%%%%%%%%%%%%%%
\section{Distance, weight, and rate}

\begin{definition}[$q$-ary entropy]
\lean{ErrorCorrectingCodes.Codeword.qaryEntropy}
\leanok
\[
H_q(p) \;=\; p\log_q(q-1) - p\log_q p - (1-p)\log_q(1-p).
\]
\end{definition}

\begin{definition}[Hamming distance]
\lean{ErrorCorrectingCodes.Codeword.hamming_distance}
\leanok
Hamming distance between two codewords, wrapping Mathlib's \lean{hammingDist}.
\end{definition}

\begin{definition}[Distance predicate for a code]
\lean{ErrorCorrectingCodes.Codeword.distance}
\leanok
\lean{distance C d} holds when $d$ is the minimum Hamming distance of $C$: there
exist distinct codewords at distance exactly $d$, and no two distinct codewords
are closer than $d$.
\end{definition}

\begin{definition}[Rate]
\lean{ErrorCorrectingCodes.Codeword.rate}
\leanok
\[
R(C) = \frac{\log|C|}{n\log|\alpha|}.
\]
\end{definition}

\begin{definition}[Weight]
\lean{ErrorCorrectingCodes.Codeword.weight}
\leanok
The Hamming weight of $c$ is $d(c, \mathbf{0})$.
\end{definition}

\begin{definition}[Maximal size at distance]
\lean{ErrorCorrectingCodes.Codeword.max_size}
\leanok
\lean{max_size n d A} asserts existence of a code of size $A$ and minimum distance $d$
that is maximal among all such codes.
\end{definition}

\begin{definition}[Hamming ball]
\lean{ErrorCorrectingCodes.Codeword.hamming_ball}
\leanok
The Hamming ball of radius $l$ around $c$:
\[
B_l(c) = \{c' : d(c',c) \le l\},
\]
implemented as a \lean{Finset}.
\end{definition}

%%%%%%%%%%%%%%%%%%%%%%%%%%%%%%%%%%%%%%%%%%%%%%%%%%%%%%%%%%%%%%%%%%%%%%%%%%%%%%%
\section{Basic lemma: distance vs.\ block length}

\begin{lemma}[Distance is at most block length]
\lean{ErrorCorrectingCodes.Codeword.dist_le_length}
\leanok
If a code $C$ has minimum distance $d$, then $d \le n$.
\end{lemma}
